\documentclass[11pt]{article}

\usepackage[T1]{fontenc}
\usepackage{lmodern}
\usepackage{geometry}
\usepackage{hyperref}
\usepackage{xcolor}
\usepackage{listings}
\usepackage{longtable}
\usepackage{booktabs}
\usepackage{multicol}
\usepackage{titlesec}

\geometry{margin=1in}

\definecolor{codebg}{RGB}{245,245,245}
\lstset{
  basicstyle=\ttfamily\small,
  backgroundcolor=\color{codebg},
  frame=single,
  columns=fullflexible,
  keepspaces=true
}

\titleformat{\section}{\Large\bfseries}{}{0em}{}
\titleformat{\subsection}{\large\bfseries}{}{0em}{}
\titleformat{\subsubsection}{\normalsize\bfseries}{}{0em}{}

\title{Bash Parameter Expansion Cheat Sheet (Full Reference)}
\author{}
\date{}

\begin{document}
\maketitle
\tableofcontents
\newpage

\section{Why this exists}
Bash parameter expansion does string manipulation inside the shell. It avoids process creation overhead from utilities such as
\texttt{basename}, \texttt{dirname}, \texttt{cut}, \texttt{sed}, \texttt{awk}, and \texttt{tr}.
In tight loops, replacing external commands with expansions can reduce runtime dramatically.

\section{Rules of thumb}
\begin{itemize}
\item Parameter expansion operates on variables, not streaming input.
\item Patterns are shell globs, not regular expressions.
\item Quote expansions unless you explicitly want word splitting or globbing.
\item Use expansions for simple delimiter and suffix transforms. Use \texttt{awk}/\texttt{sed} for regex-heavy parsing.
\end{itemize}

\section{Core operators for trimming}
\subsection{Prefix and suffix removal}
\begin{longtable}{lll}
\toprule
Form & Meaning & Match type \\
\midrule
\texttt{\$\{v\#pattern\}}   & remove from start & shortest \\
\texttt{\$\{v\#\#pattern\}} & remove from start & longest \\
\texttt{\$\{v\%pattern\}}   & remove from end   & shortest \\
\texttt{\$\{v\%\%pattern\}} & remove from end   & longest \\
\bottomrule
\end{longtable}

\subsection{Mental model}
\begin{itemize}
\item \texttt{\#} trims from the left.
\item \texttt{\%} trims from the right.
\item doubling means greedy.
\item patterns are globs such as \texttt{*}, \texttt{?}, and \texttt{[abc]}.
\end{itemize}

\subsection{Examples}
\begin{lstlisting}
v="aa/bb/cc.txt"

${v#*/}    -> bb/cc.txt
${v##*/}   -> cc.txt
${v%/*}    -> aa/bb
${v%%/*}   -> aa
\end{lstlisting}

\section{Path munging cookbook}
Assume:
\begin{lstlisting}
p="/var/log/nginx/access.log"
\end{lstlisting}

\subsection{Basename and dirname}
\begin{lstlisting}
base=${p##*/}   # access.log
dir=${p%/*}     # /var/log/nginx
\end{lstlisting}

Utility equivalents:
\begin{lstlisting}
base=$(basename -- "$p")
dir=$(dirname -- "$p")
\end{lstlisting}

\subsection{Extension and stem}
\begin{lstlisting}
stem=${base%.*}   # access
ext=${base##*.}   # log
\end{lstlisting}

For multi-extensions:
\begin{lstlisting}
f="archive.tar.gz"
stem1=${f%.*}     # archive.tar
ext1=${f##*.}     # gz
stem2=${f%%.*}    # archive
\end{lstlisting}

\subsection{Parent directory name (last segment of dirname)}
\begin{lstlisting}
parent=${dir##*/}   # nginx
\end{lstlisting}

\subsection{Join, normalize, and guard}
\begin{lstlisting}
# join pieces safely
out="${dir}/${stem}.csv"

# ensure non-empty before use (fails fast)
: "${p:?path required}"

# default if empty
: "${dir:=.}"
\end{lstlisting}

\subsection{Strip a known prefix if present}
\begin{lstlisting}
x="/mnt/data/file.txt"
y=${x#/mnt/}     # data/file.txt
\end{lstlisting}

\subsection{Strip a known suffix if present}
\begin{lstlisting}
x="report.csv.bak"
y=${x%.bak}      # report.csv
\end{lstlisting}

\subsection{Remove everything up to the first slash}
\begin{lstlisting}
x="a/b/c"
y=${x#*/}        # b/c
\end{lstlisting}

\subsection{Remove everything after the last slash}
\begin{lstlisting}
x="a/b/c"
y=${x%/*}        # a/b
\end{lstlisting}

\subsection{Split around delimiters, cut-style}
\begin{lstlisting}
v="host01.domain.com"
short=${v%%.*}   # host01
rest=${v#*.}     # domain.com
\end{lstlisting}

\section{Substring extraction}
\subsection{By position}
\begin{lstlisting}
s="abcdef"

${s:2}     -> cdef
${s:2:3}   -> cde
${s: -2}   -> ef   # note the space before -2
\end{lstlisting}

\section{Search and replace}
\begin{longtable}{ll}
\toprule
Form & Meaning \\
\midrule
\texttt{\$\{v/pat/repl\}}    & replace first match \\
\texttt{\$\{v//pat/repl\}}   & replace all matches \\
\texttt{\$\{v/\#pat/repl\}}  & replace at start \\
\texttt{\$\{v/\%pat/repl\}}  & replace at end \\
\bottomrule
\end{longtable}

\subsection{Examples}
\begin{lstlisting}
v="a-b-c"
${v/-/_}     -> a_b-c
${v//-/_}    -> a_b_c

p="/a/b/c.txt"
${p//\//_}   -> _a_b_c.txt
\end{lstlisting}

\subsection{Compare to sed}
\begin{lstlisting}
sed 's/-/_/'     # first
sed 's/-/_/g'    # all
\end{lstlisting}

Use \texttt{sed} when you need regex capture groups or complex patterns.

\section{Case conversion (bash 4+)}
\begin{lstlisting}
v="MiXeD"
${v,,}   -> mixed
${v^^}   -> MIXED
${v,}    -> miXeD
${v^}    -> MiXeD
\end{lstlisting}

\section{Defaults, errors, conditionals}
\begin{longtable}{ll}
\toprule
Form & Meaning \\
\midrule
\texttt{\$\{v:-d\}}   & use d if unset or empty, do not assign \\
\texttt{\$\{v:=d\}}   & assign d if unset or empty \\
\texttt{\$\{v:?msg\}} & error if unset or empty \\
\texttt{\$\{v:+alt\}} & use alt if set and non-empty \\
\bottomrule
\end{longtable}

\section{Length}
\begin{lstlisting}
len=${#v}
\end{lstlisting}

\section{Whitespace trimming (verbose, glob-based)}
\begin{lstlisting}
# leading
v="${v#"${v%%[!$' \t\r\n']*}"}"
# trailing
v="${v%"${v##*[!$' \t\r\n']}"}"
\end{lstlisting}

\section{When to use what}
\begin{longtable}{p{0.48\textwidth}p{0.48\textwidth}}
\toprule
Prefer parameter expansion & Prefer external tools \\
\midrule
tight loops, thousands of iterations & complex regex parsing \\
simple delimiters, suffix/prefix logic & multi-field logic, arithmetic \\
path manipulation and renaming & streaming text pipelines \\
\bottomrule
\end{longtable}

\end{document}
